% =====================================================================
% SEC_09 — Considerações Finais e Trabalho Futuro
% =====================================================================

\section*{9. Considerações Finais e Trabalho Futuro}
\addcontentsline{toc}{section}{9. Considerações Finais e Trabalho Futuro}

\nivel{0} Chegamos ao fim. Esta seção amarra o fio narrativo: retomamos o problema
de pesquisa da Seção~2, verificamos como cada seção contribuiu para resolvê-lo e
definimos um \emph{roadmap concreto} de trabalho futuro, com etapas verificáveis e
métricas de aceitação. Nada de promessas vagas: cada item tem critério e
responsável (o protocolo TDD do ecossistema).

\subsection*{9.1 Síntese: o caminho percorrido}

\nivel{1} O manual seguiu uma sequência lógica que pode ser resumida em um único
fluxo de raciocínio:

\begin{center}
\begin{tabularx}{\textwidth}{lX}
\toprule
\textbf{Etapa} & \textbf{O que se estabeleceu} \\
\midrule
Seção 1 & Leitor é orientado sobre como ler (trilha N0--N2). \\
Seção 2 & Problema de pesquisa (PP), 4 RQs e hipótese H1 são formalizados. \\
Seção 3 & Base teórica: BM25, RAG, raciocínio, MMR e Recamán. \\
Seção 4 & Estado atual: pipeline implementado e \emph{lacuna} de diversidade. \\
Seção 5 & Proposta pós-Recamán: 4 artefatos e posicionamento aditivo. \\
Seção 6 & Cálculos: métricas, custo $O(N)$ e conexão com "lost in the middle". \\
Seção 7 & Implementações: algoritmo, interface e reprodutibilidade. \\
Seção 8 & Mapas: transversalidade e impacto sem quebrar interfaces. \\
\bottomrule
\end{tabularx}
\end{center}

\begin{pontochave}
O fio condutor único: \emph{um problema real} (concentração de resultados)
$\rightarrow$ \emph{uma hipótese} (diversificação determinística via Recamán)
$\rightarrow$ \emph{uma especificação} (artefatos + métricas) $\rightarrow$ \emph{uma
implementação reproduzível}.
\end{pontochave}

\subsection*{9.2 Retomada do problema e das hipóteses}

\nivel{2} Voltando à Seção~2: o problema de pesquisa (PP) pedia uma diversificação
estruturada determinística, pós-ranqueamento, com custo controlado. O manual
\emph{responde} a cada componente:

\begin{itemize}
  \item \textbf{Determinismo} --- respondido pela sequência de Recamán
        (Eq.~\ref{eq:recaman}), que é uma função matemática canônica, sem
        \emph{seed} (Seções 5 e 7).
  \item \textbf{Pós-ranqueamento} --- o \texttt{RecamanDiversifier} opera sobre o
        ranking já ordenado, sem alterar o ranqueador (Seção 5).
  \item \textbf{Custo controlado} --- demonstrado em $O(N)$ (Eq.~\ref{eq:cost}),
        não dominante no pipeline (Seção 6).
\end{itemize}

\noindent A hipótese H1, portanto, permanece \textbf{plausível e mensurável}, mas
\textbf{não verificada}: sua confirmação exige implementação e experimentação
empírica. Este manual a \emph{formula} e a \emph{prepara}, sem a \emph{provar}.

\subsection*{9.3 Roadmap de trabalho futuro (concreto e verificável)}

\nivel{1} Definimos um roadmap em fases, cada uma com \textbf{critério de aceitação
mensurável} e alinhado às RQs. Este é o plano \emph{ação-ação}, passível de virar
specs SDD/TDD:

\begin{enumerate}[label=\textbf{Fase \arabic*.}]
  \item \textbf{Implementação do \texttt{RecamanDiversifier}} (RQ2, RQ3).
        \begin{itemize}
          \item \emph{Ação:} criar \texttt{rag/recaman.py} com a interface da
                Seção~7, seguindo TDD (RED$\to$GREEN$\to$REFACTOR).
          \item \emph{Critério:} testes deterministas ($a_0\!..\!a_5 =
                [0,1,3,6,2,7]$), terminação para todos os $N$,$k$, e
                $\text{Div}(\mathcal{S})$ mensurável.
        \end{itemize}
  \item \textbf{Métrica de diversidade no \texttt{EnhancedRAG.metrics()}} (RQ1).
        \begin{itemize}
          \item \emph{Ação:} adicionar o campo $\text{Div}(\mathcal{S})$
                (Eq.~\ref{eq:diversity}) ao painel de métricas existente, ligado a
                \texttt{AnchorResolver}.
          \item \emph{Critério:} painel passa a reportar diversidade; teste
                garante $\text{Div}\in[0,1]$ e coerência com \emph{groundedness}.
        \end{itemize}
  \item \textbf{Experimento de coorte em corpus-piloto} (RQ1, RQ4).
        \begin{itemize}
          \item \emph{Ação:} comparar diversidade e groundedness entre estado
                atual, MMR e Recamán em um corpus real do ecossistema.
          \item \emph{Critério:} relatório com $\text{Div}$, groundedness e
                \emph{baseline} justo (mesmos dados, mesmas seeds, mesmo oráculo).
        \end{itemize}
  \item \textbf{Teste de regressão de interface} (RQ4).
        \begin{itemize}
          \item \emph{Ação:} garantir que fluxos (acadêmico, jurídico, clínico,
                MIRA) continuem consumindo o RAG sem mudança de contrato.
          \item \emph{Critério:} suíte de regressão verde sobre a interface pública
                dos fluxos.
        \end{itemize}
  \item \textbf{Registro evolutivo} (governança).
        \begin{itemize}
          \item \emph{Ação:} registrar cada ciclo no \texttt{evolution/cycles.json}
                e no MetaBus, conforme o protocolo do ecossistema.
          \item \emph{Critério:} ciclo R456+ registrado com score e lições;
                MetaBus acessível.
        \end{itemize}
\end{enumerate}

\begin{intuicao}
O roadmap é como a "lista de tarefas de um projeto": cada linha tem um \emph{o
que fazer} (ação), um \emph{como saber que deu certo} (critério) e uma \emph{ligação
com a pergunta de pesquisa} (RQ). Isso transforma intenções em trabalho auditável.
\end{intuicao}

\subsection*{9.4 Contribuições esperadas (com rigor)}

\nivel{2} As contribuições esperadas deste trabalho, de forma \emph{cautelosa} e
sem \emph{overclaim}, são:

\begin{enumerate}[label=\arabic*.]
  \item \textbf{Contribuição metodológica:} uma ponte determinística entre teoria
        dos números (Recamán) e diversificação de recuperação, oferecendo um
        \emph{baseline reproduzível} e comparável (Seções 5--6).
  \item \textbf{Contribuição de engenharia:} uma especificação aditiva e de baixo
        acoplamento para o ecossistema, com custo $O(N)$ e preservação de
        interface (Seções 6 e 8).
  \item \textbf{Contribuição de medição:} o fechamento da lacuna de \emph{métrica de
        diversidade} do painel atual (Seção 4), permitindo avaliar a hipótese H1.
\end{enumerate}

\begin{pontochave}
Diferentemente de relatos que superestimam resultados, este manual declara
\emph{contribuições esperadas} e \emph{hipotéticas}, vinculadas a RQs e métricas
— nunca como fatos consumados.
\end{pontochave}

\subsection*{9.5 Considerações éticas e de rigor final}

\nivel{2} Reafirmamos os compromissos éticos e de rigor que atravessaram todo o
documento:

\begin{itemize}
  \item \textbf{Anti-overclaim:} nenhuma certificação externa, segurança absoluta
        ou disponibilidade universal é alegada.
  \item \textbf{Transparência:} estado atual e proposta estão separados e
        rotulados; nenhum componente proposto é apresentado como implementado.
  \item \textbf{Mensurabilidade:} todo impacto potencial está vinculado a uma RQ e
        a uma métrica (Eq.~\ref{eq:diversity}), a ser validada empiricamente.
  \item \textbf{Reprodutibilidade:} algoritmos, métricas e até a compilação deste
        manual são reproduzíveis (Seção 7).
\end{itemize}

\begin{pontochave}
O fecho deste manual não é um ponto final, mas um \emph{ponto de partida}:
deixamos o problema formulado, a hipótese preparada e o roadmap aberto, prontos
para migrar da \emph{especificação} para a \emph{implementação verificável}.
\end{pontochave}
